%! From Combinations to Solutions --- Setting Up Ax = b — Unit 2, Lesson 2.0 — Slides (Beamer)
\documentclass[aspectratio=169, 11pt]{beamer}
\usepackage{linalg-beamer}   % provides \burgundyheader{} and \sectionlabel[color]{}
\newcommand{\xx}{\mathbf{x}}
\newcommand{\bb}{\mathbf{b}}

\begin{document}

% TITLE SLIDE (CourseName is not defined in beamer, so the name is literal)
{
\setbeamercolor{background canvas}{bg=burgundy}
\begin{frame}[plain]
  \vspace{2.8cm}\hspace{0.55cm}%
  \begin{minipage}{0.9\textwidth}
    {\color{goldacc}\bfseries\small LESSON 2.0}\\[0.5cm]
    {\color{white}\bfseries\LARGE From Combinations to Solutions}\\[0.3cm]
    {\color{white}\bfseries\large Setting Up $A\xx=\bb$}\\[1.0cm]
    {\color{blushmid}\small Linear Algebra~$\cdot$~Unit 2: Solving Linear Equations $A\xx=\bb$}
  \end{minipage}
\end{frame}
}

% HOOK
\begin{frame}[t, plain]
  \burgundyheader{Hit the target}
  \sectionlabel{Hook}
  \begin{columns}[T]
    \begin{column}{0.60\textwidth}
      One scoop of \textbf{Base A} $=1$ cup nuts, $3$ cups raisins.\\[4pt]
      One scoop of \textbf{Base B} $=2$ cups nuts, $2$ cups raisins.\\[8pt]
      A recipe needs \textbf{$5$ cups nuts} and \textbf{$11$ cups raisins}.\\[10pt]
      \textbf{Question.} How many scoops of each --- and is there only one answer?
    \end{column}
    \begin{column}{0.38\textwidth}
      \centering
      \[
        x_1\begin{bmatrix} 1 \\ 3 \end{bmatrix}
        + x_2\begin{bmatrix} 2 \\ 2 \end{bmatrix}
        = \begin{bmatrix} 5 \\ 11 \end{bmatrix}
      \]
    \end{column}
  \end{columns}
\end{frame}

% RUNNING UNIT 1 BACKWARDS
\begin{frame}[t, plain]
  \burgundyheader{Running Unit 1 backwards}
  \sectionlabel{Set-up}
  \begin{itemize}
    \setlength{\itemsep}{6pt}
    \item In Lesson~1.3 we \textbf{built} the target: given $A$ and $\xx$, compute $\bb=A\xx$.
    \item Now the target $\bb$ is \textbf{given} --- we hunt for the weights $\xx$.
  \end{itemize}
  \vspace{4pt}
  \[
    A = \begin{bmatrix} 1 & 2 \\ 3 & 2 \end{bmatrix}, \qquad
    \bb = \begin{bmatrix} 5 \\ 11 \end{bmatrix}, \qquad
    \text{solve } A\xx=\bb.
  \]
  \begin{block}{To solve $A\xx=\bb$}
    Find the weights $\xx$ that combine the \textbf{columns} of $A$ into the target $\bb$.
  \end{block}
\end{frame}

% TWO LANGUAGES
\begin{frame}[t, plain]
  \burgundyheader{One problem, two languages}
  \sectionlabel[goldacc]{Columns vs.\ Rows}
  \begin{columns}[T]
    \begin{column}{0.5\textwidth}
      \textbf{Column view} (down the columns):
      \[
        x_1\begin{bmatrix} 1 \\ 3 \end{bmatrix}
        + x_2\begin{bmatrix} 2 \\ 2 \end{bmatrix}
        = \begin{bmatrix} 5 \\ 11 \end{bmatrix}
      \]
    \end{column}
    \begin{column}{0.5\textwidth}
      \textbf{Row view} (across the rows):
      \[
        \begin{aligned}
          x_1 + 2x_2 &= 5 \\
          3x_1 + 2x_2 &= 11
        \end{aligned}
      \]
    \end{column}
  \end{columns}
  \vspace{8pt}
  \begin{block}{Same numbers, two readings}
    Down the columns is a \textbf{combination}; across the rows is a \textbf{system of equations}.
  \end{block}
\end{frame}

% ELIMINATION
\begin{frame}[t, plain]
  \burgundyheader{Solve by elimination}
  \sectionlabel{The key move}
  Both equations carry a $2x_2$, so \textbf{subtract} the first from the second:
  \[
    (3x_1+2x_2)-(x_1+2x_2) = 11-5
    \quad\Longrightarrow\quad 2x_1 = 6 \quad\Longrightarrow\quad x_1 = 3.
  \]
  Back-substitute into $x_1+2x_2=5$: \quad $x_2 = 1$.
  \begin{block}{Check by rebuilding $\bb$}
    $3\begin{bsmallmatrix} 1 \\ 3 \end{bsmallmatrix}
     + 1\begin{bsmallmatrix} 2 \\ 2 \end{bsmallmatrix}
     = \begin{bsmallmatrix} 5 \\ 11 \end{bsmallmatrix}$ \checkmark\quad
    $3$ scoops of A, $1$ of B. \textbf{Subtracting to kill a variable $=$ elimination.}
  \end{block}
\end{frame}

% ONE / NONE / MANY
\begin{frame}[t, plain]
  \burgundyheader{One, none, or infinitely many}
  \sectionlabel[goldacc]{Geometry}
  Each equation is a \textbf{line}; a solution is where they meet.
  \begin{itemize}
    \setlength{\itemsep}{5pt}
    \item Lines \textbf{cross once} $\Rightarrow$ \textbf{one} solution (consistent).
    \item Lines \textbf{parallel} $\Rightarrow$ \textbf{no} solution (inconsistent).
    \item \textbf{Same line} $\Rightarrow$ \textbf{infinitely many} solutions.
  \end{itemize}
  \vspace{4pt}
  \begin{block}{Back to reachability (Lesson 1.3)}
    $A\xx=\bb$ has a solution exactly when $\bb$ is \textbf{reachable} --- in the column space of $A$.
  \end{block}
\end{frame}

% ROADMAP / PREVIEW
\begin{frame}[t, plain]
  \burgundyheader{The road through Unit 2}
  \sectionlabel{Connections}
  \textbf{Today:} $A\xx=\bb$ as columns and as equations; solve $2\times 2$ by elimination.\\[8pt]
  \begin{itemize}
    \setlength{\itemsep}{4pt}
    \item \textbf{2.1} The Idea of Elimination --- pivots, done systematically.
    \item \textbf{2.2} Elimination Matrices and Inverse Matrices.
    \item \textbf{2.3} Matrix Computations and $A=LU$.
    \item \textbf{2.4} Permutations and Transposes.
  \end{itemize}
\end{frame}

\end{document}
